\documentclass[openright，twoside]{book}

\usepackage{tjustyle}

% 根据需求导入宏包，自定义命令。以下内容如不需要都可删除。

\RequirePackage{threeparttable} % 用于给表格添加注释
\RequirePackage{threeparttablex} % 用于给长表格添加注释
\RequirePackage{longtable} % 用于绘制长表格

% 定义表格列宽的命令
\newcolumntype{L}[1]{>{\raggedright\let\newline\\\arraybackslash\hspace{0pt}}m{#1}}
\newcolumntype{C}[1]{>{\centering\let\newline\\\arraybackslash\hspace{0pt}}m{#1}}
\newcolumntype{R}[1]{>{\raggedleft\let\newline\\\arraybackslash\hspace{0pt}}m{#1}}


\graphicspath{{figures/}} % 指定图片保存地址

\begin{document}

% 扉页，可根据实际情况修改`\vsapce'值以调整各部分的垂直距离，使内容保持在一页之内

\thispagestyle{empty} % 去除页眉页脚
\begin{center} % 本页内容居中

\vspace*{0cm}

\zihao{2}\textbf{中文题目}

\vspace{0.8cm}

\zihao{4}\textbf{English Title}

\vspace{4cm}

\zihao{-4} % 切换为小四字号绘制下面两个表格

% 基本信息表
\begin{table}[h] 
    \centering 
    \renewcommand*{\arraystretch}{1.35} % 设置行高
    \begin{tabular}{cc}
        \textbf{一级学科:} & \textbf{\underline{\makebox[14em][c]{在此输入}}} \\
        \textbf{研究方向:} & \textbf{\underline{\makebox[14em][c]{输入于此}}} \\
        \textbf{作者姓名:} & \textbf{\underline{\makebox[14em][c]{某某某}}} \\
        \textbf{指导教师:} & \textbf{\underline{\makebox[14em][c]{某某某}}} \\
    \end{tabular}
\end{table}

% 答辩信息表，可根据实际答辩委员数增加行数，应修改 \multirow 的第一个参数
\begin{table}[h] 
    \centering 
    \renewcommand*{\arraystretch}{1.35} % 设置行高
    \begin{tabularx}{\textwidth}{|*{4}{>{\centering\arraybackslash}X|}}
    \hline
    \textbf{答辩日期}               & \multicolumn{3}{c|}{2022年11月29日}               \\ \hline
    \textbf{答辩委员会}             & \textbf{姓名}  & \textbf{职称} & \textbf{工作单位}  \\ \hline
    \textbf{主席}                  &               &               &                  \\ \hline
    \multirow{6}{*}{\textbf{委员}} &               &               &                  \\ \cline{2-4} 
                                  &               &               &                  \\ \cline{2-4}
                                  &               &               &                  \\ \cline{2-4}
                                  &               &               &                  \\ \cline{2-4}
                                  &               &               &                  \\ \cline{2-4}
                                  &               &               &                  \\ \hline
    \end{tabularx}
\end{table}

\vspace{0cm}

\zihao{4}{天津大学某某学院\\二〇二二年十一月}

\end{center} % 停止居中
 % 扉页
% 声明授权页

\clearpage{\cleardoublepage}
\thispagestyle{empty}

\vspace*{1cm}

\begin{center}
    \zihao{-2}{独创性声明}    
\end{center}

\vspace{0.5cm}

\zihao{-4}
本人声明所呈交的学位论文是本人在导师指导下进行的研究工作和取得的研究成果，除了文中特别加以标注和致谢之处外，论文中不包含其他人已经发表或撰写过的研究成果，也不包含为获得{\underline{\zihao{4}\textbf{\kaishu{~~天津大学~~}}}}或其他教育机构的学位或证书而使用过的材料。与我一同工作的同志对本研究所做的任何贡献均已在论文中作了明确的说明并表示了谢意。

\vspace{1.5cm} 

\noindent 学位论文作者签名：\makebox[4cm][s]{} 签字日期：\makebox[1.5cm][s]{}年\makebox[1cm][s]{}月\makebox[1cm][s]{}日

\vspace{2.5cm}

\begin{center}
    \zihao{-2}{学位论文版权使用授权书}    
\end{center}

\vspace{1cm}

\zihao{-4}
本学位论文作者完全了解{\underline{\zihao{4}{\textbf{\kaishu{~~天津大学~~}}}}}有关保留、使用学位论文的规定。特授权{\underline{\zihao{4}{\textbf{\kaishu{~~天津大学~~}}}}}可以将学位论文的全部或部分内容编入有关数据库进行检索，并采用影印、缩印或扫描等复制手段保存、汇编以供查阅和借阅。同意学校向国家有关部门或机构送交论文的复印件和磁盘。

（保密的学位论文在解密后适用本授权说明）

\vspace{2.2cm}

\noindent 学位论文作者签名：\makebox[4cm][s]{} 导师签名：\makebox[4cm][s]{}

\vspace{0.8cm}

\noindent 签字日期：\makebox[1.5cm][s]{}年\makebox[1cm][s]{}月\makebox[1cm][s]{}日 
\qquad 签字日期：\makebox[1.5cm][s]{}年\makebox[1cm][s]{}月\makebox[1cm][s]{}日
 % 独创性声明

\frontmatter % 开始有页码（大写罗马）

% 中文摘要页

\chapter*{\songti\zihao{2}\textbf{摘 \qquad 要}}
\thispagestyle{noheader}

\zhlipsum[1-4] 

\bigskip

\noindent {\zihao{4}\textbf{关键词：}} 
关键词，关键词，关键词，关键词，关键词

% 英文摘要页

\chapter*{\songti\zihao{2}\textbf{ABSTRACT}}
\thispagestyle{noheader}

\lipsum[1-5] 

\bigskip

\noindent {\zihao{4}\textbf{KEY WORDS:} }
Keyword, Keyword, Keyword, Keyword, Keyword
 % 中英文摘要

\tableofcontents % 目录

\mainmatter % 进入正文部分，页码变为阿拉伯数字


\chapter{章标题（一级标题）}

\section{节标题（二级标题）}

这里测试大段文字：

学位论文是学位授予单位研究生培养质量的重要标志，是研究生本人在学期间从事学术研究的成果体现，是对研究生综合能力的考核，也是研究生申请学位的主要依据。为保证我校研究生学位论文的质量，实现我校学位论文格式的规范化，参照GB/T 7713.1—2006《学位论文编写规则》，GB 7713-87《科学技术报告、学位论文和学术论文的编写格式》和GB/T 7714—2005《文后参考文献著录规则》等标准，特制定《天津大学关于博士、硕士学位论文统一格式的规定》。

\begin{enumerate}[label={（\arabic*）}]
    \item 测试列表条目
    \item 测试列表条目 
    \begin{enumerate}[label={（\roman*）}]
        \item 嵌套列表条目
        \item 嵌套列表条目
    \end{enumerate}
\end{enumerate}

这里测试插入图片：
\begin{figure}[!h]
    \centering\small
    \includegraphics[width=3cm]{tju_logo.png}
    \caption{天津大学校徽}
    \label{fig:logo}
\end{figure}
图~\ref{fig:logo} 所示为一张图片~\cite{mittelbach2004}。

\subsection{小节标题（三级标题）}

\chapter{另一个章标题（一级标题）}

测试新一章是否正确地从奇数页开始。 % 导入正文示例

\bibliographystyle{gbt7714} % 指定参考文献格式文件
\bibliography{references} % 指定参考文献数据文件

\appendix
\chapter{附录标题}

附录内容。 % 添加附录

\backmatter % 以后的章节没有编号

\chapter{发表论文和参加科研情况说明}

发表论文和参加科研情况说明。
\begin{enumerate}[label={[\arabic*]}]
    \item 我的一篇论文
    \item 我的一篇论文
    \item 我的一篇论文
\end{enumerate} % 发表论文和参加科研情况说明
\chapter{致\qquad 谢}

致谢内容。 % 致谢

\end{document}
